As AI models become smarter, how do we ensure that they continue to act in the interest of humans? A long line of work~\citep{} has argued that this is a hard problem, and recent events such as the infamous HuggingFace incident~\citep{} have demonstrated the difficulties in practice. 

One way to study this problem more formally has been termed \emph{scalable oversight}. Here, a smarter-than-human AI is typically modelled as an algorithm that can solve computational problems in some large complexity class, and a human is modelled as an algorithm in a much smaller complexity class. It is assumed that given a context $x$, computing a human-aligned action $y$ is hard for the human but possible for the AI. How can the human verify whether the AI is, in fact, calculating $y$, as opposed to some other bad action $y'$? Of course, without any extra assumptions this is impossible. The scalable oversight literature assumes that the AI sits within some level of the polynomial hierarchy, thus providing a means to \emph{certify} whether AI's solution is correct. For example, for the first level of the hierarchy, one assumes the AI to be in $\operatorname{NP} \cap \operatorname{coNP}$ and the human in P, thus implying that there exists a certificate $c$ that the human can look at to verify whether $y$ is the correct answer to $x$. For the second level, this implies there exists a certificate $c_1$ such that for all certificates $c_2$, the human's response to seeing $x, y, c_1$, and $c_2$ is to judge $y$ to be the correct answer to $x$. For the $k^\text{th}$ level, this looks like debate where two debaters alternately produce certificates in a way that the first one has a winning strategy if and only if $y$ is the correct answer. 
